%% Discussion. Numbers: tables/numbers.tex.
\paragraph{What the independent test confirmed.} A feasibility classifier trained on household scans transferred to evolved shapes it had never seen: the learned order examined fewer candidates than the original planner (P1), and adding the robustness and centring terms cost nothing in efficiency relative to the plain learned order (P3). The gain was smaller than in development (\confPoneRatio{} vs \devPoneRatio), because the original planner's heuristic order is already good on gripper-scaled EGAD shapes. The population-level view matters here: \confNonViable{} of \confObjects{} objects had no viable candidate at all, and no ordering can change that.

\paragraph{What it did not confirm.} The robustness advantage that motivated FORTIS-R Joint in development (\devJointMzeroRobust{} points) did not replicate; on EGAD Joint was non-inferior to the original planner only at the pre-registered friction value, narrowly, and fell below it at the other torsional-friction values. Load tolerance was lower than the original planner's. The historical claims that a centring term improves load tolerance (H1) and that a learned robustness model improves robustness (H2) also failed to replicate. Two explanations are consistent with the data, and the study cannot separate them: the robustness model was trained on development labels from a different population (household scans, the previous rise-and-slip criterion), and the original planner's grasps are already highly robust on EGAD (\confRobustMzero\%), leaving little room. Either way, a robustness term learned on one object population should not be assumed to help on another.

\paragraph{Physics modelling decides load-tolerance conclusions.} Two choices that are rarely reported changed conclusions: the torsional friction coefficient (between \devFlipHigh\% and \devFlipLow\% of verdicts changed over the tested range) and whether the gripper floats or hangs from an articulated arm (load contrasts reversed between the two in development). Moving load tolerance out of the primary family before the test was justified by these development results; the confirmation then agreed with the articulated direction. Claims about load tolerance from a single floating-gripper simulation should be treated with caution.

\paragraph{Efficiency in candidates is not efficiency in seconds.} For this planner, generating the catalogue costs several times more than searching it, and the learned order only shortens the search. With the present two-process implementation the learned orders were roughly twice as slow end to end (H4 failed). Integrating the prior into the planner removes the duplicated generation; our estimate of that implementation is close to parity with the original planner, not faster. The number of examined candidates remains the right measure when candidate evaluation is expensive (e.g. with full motion planning or physics in the loop), but the claim of a wall-clock saving does not hold here.

\paragraph{Foundation-model features.} Frozen DINOv2, DINOv3 and OpenCLIP global descriptors, DINOv3 patch tokens at the contacts, and a listwise ranking loss did not improve on a gradient-boosting model over 19 hand-designed candidate features, in development or confirmation, once permuted and untrained controls were considered. With the learned order the first examined candidate is already viable for most objects, so there is little left for richer object descriptors---or for few-shot adaptation---to gain. This is a bounded negative result: it concerns frozen descriptors of grey renders for ordering a planner's candidates, not visual grasp learning in general.

\paragraph{Practical guidance.} For a complete sampling-based planner: (i) learn the examination order from exhaustively labelled catalogues and evaluate it with an exact replay of the planner; (ii) run it in round robin with the original order to bound the worst case at twice the original cost; (iii) validate dynamic claims across friction and simulator-fidelity variations and on an independent population before relying on a learned robustness term; (iv) measure end-to-end time, not only examined candidates.

\paragraph{Limitations.} Everything is simulated, with assumed masses, uniform density and uncalibrated friction; no real robot was used. One cell and one 80~mm gripper were confirmed; the 40 and 120~mm results are development-only stroke variants. EGAD shapes are gripper-scaled evolved objects, not household items, and their collision geometry is a convex decomposition. The articulated model uses the cell's link inertias and an ideal computed-torque controller without arm collisions. No learned external grasp detector was reproduced (GPD would need a C++ adapter; Contact-GraspNet a GPU; Dex-Net a legacy software stack); the only external comparison is a reimplemented analytic criterion. Two secondary hypotheses (H2 in particular) were underpowered, and non-significant results there are not evidence of equivalence.
